\title{xLSTM: Extended Long Short-Term Memory}

\begin{abstract}
The introduction of the constant error carousel and gating mechanisms in the 1990s laid the foundation for the development of the Long Short-Term Memory (LSTM). Over the years, LSTMs have demonstrated remarkable robustness and played a pivotal role in the evolution of deep learning, notably forming the basis for the earliest Large Language Models (LLMs). Despite these achievements, the emergence of Transformer architectures, built upon highly parallelizable self-attention mechanisms, signified a paradigm shift by surpassing LSTM performance at large scale. This leads us to investigate a straightforward question: What are the capabilities of LSTMs in language modeling when expanded to billions of parameters, employing modern LLM advancements and addressing LSTM-specific shortcomings? Our work begins by proposing exponential gating enhanced with dedicated normalization and stabilization methods. Furthermore, we present alterations to the canonical LSTM memory system, resulting in: (i) sLSTM, which features scalar memory and update, coupled with a novel memory mixing scheme, and (ii) mLSTM, which achieves full parallelism by employing a matrix memory and a covariance-based update mechanism. By incorporating these extended LSTM variants into residual block frameworks, we develop xLSTM blocks, which can be stacked to build xLSTM models. The combined effect of exponential gating and revised memory designs significantly strengthens the xLSTM’s ability to compete with contemporary Transformers and State Space Models, demonstrating promising outcomes in both effectiveness and scalability.\\
Code available at: \url{https://github.com/NX-AI/xlstm}
\end{abstract}

\section{Introduction}
\label{sec:intro}

Long Short-Term Memory (LSTM) mechanisms~\citep{Hochreiter:91,Hochreiter:97nips,Hochreiter:97}, specifically the constant error carousel and gating components, were proposed as solutions to address the issue of vanishing gradients in recurrent neural networks~\citep{Hochreiter:91,Hochreiter:00book}:

\begin{align}
\label{eq:lstm_idea}
  c_t \ = \ f_t \ c_{t-1} \ + \ 
          i_t \  z_t \ , \quad  
          h_t \ = \ o_t \ \psi({c_t}) \ . 
\end{align}

The constant error carousel refers to the process where the cell state $c_{t-1}$ (depicted in green) receives an additive modification via the cell inputs $z_t$, with this process regulated by the sigmoid gates (illustrated in blue). The input gate $i_t$ and the forget gate $f_t$ determine how this update unfolds, whereas the output gate $o_t$ is responsible for managing the memory cell's output, specifically the hidden state $h_t$. After normalization or squashing by $\psi$, the cell state is then modulated by the output gate to produce the hidden state.

LSTMs have seen widespread adoption across a variety of fields \citep{Hochreiter:01,Hochreiter:07,Schmidhuber:15}, maintaining dominance in text generation tasks prior to the emergence of Transformers in 2017~\citep{Vaswani:17}. Their strong performance has been substantiated in multiple sequence-based applications, including but not limited to text generation~\citep{Graves:13,Karpathy:15blog}, handwriting synthesis~\citep{Graves:13}, sequence-to-sequence machine translation~\citep{Sutskever:14nips}, program evaluation~\citep{Zaremba:14arxiv}, image captioning~\citep{Karpathy:15,Hossain:19}, source code generation~\citep{Karpathy:15blog}, rainfall-runoff prediction~\citep{Kratzert:18,Kratzert:19}, and flood warning hydrological modeling~\citep{Nearing:24}. In the context of reinforcement learning, LSTMs have outperformed other sequence modeling approaches, as evidenced by their use in systems such as AlphaStar for StarCraft II~\citep{Vinyals:17short}, OpenAI Five for Dota 2~\citep{OpenAI:19}, and controllers for magnetic systems in nuclear fusion research~\citep{Degrave:22short}. Notably, LSTMs are proficient in abstraction; they effectively distill and retain semantic representations within their memory units~\citep{Karpathy:15blog}, as demonstrated by the identification of number and syntax neurons~\citep{Lakretz:19}, linguistic features~\citep{Bau:19}, and sentiment neurons~\citep{Radford:17arxiv}. Despite ongoing advancements in machine learning, LSTMs remain integral to cutting-edge applications~\citep{Degrave:22short,Nearing:24}, highlighting their enduring utility and resilience.

Although LSTMs have achieved remarkable success, they are subject to three principal limitations: (i) Difficulty in revising stored information. This issue is illustrated using the {\it Nearest Neighbor Search} task (refer to Appendix~\ref{sec:appNearestNeighborSearch}): Given a reference vector, the model must search through a sequence to identify the most similar vector and output its associated value at the sequence’s conclusion. The left side of Figure~\ref{fig:lstmProblems} presents the mean squared error for this task. Standard LSTMs have trouble updating stored values when a more similar vector appears later in the sequence, whereas our novel xLSTM addresses this deficiency via exponential gating. (ii) Restricted storage capability, necessitating the compression of information into scalar cell states. This limitation is demonstrated through the {\it Rare Token Prediction} problem. In the right side of Figure~\ref{fig:lstmProblems}, the perplexity in token prediction for Wikitext-103~\citep{Merity:17} is displayed for partitions grouped by token frequency. LSTMs display inferior performance for infrequent tokens due to their storage constraints, whereas xLSTM leverages matrix-based memory to overcome this challenge. (iii) Absence of parallelizability arising from memory interdependence: hidden-to-hidden state connections across time steps enforce strict sequential computation.

The shortcomings inherent to LSTM architectures have catalyzed the adoption of Transformers~\citep{Vaswani:17} within language modeling tasks. An open question thus arises: if the aforementioned constraints are addressed and LSTMs are scaled to match the size of today’s Large Language Models, what levels of language modeling performance become attainable?

\section{Extended Long Short-Term Memory}
\label{sec:xLSTM}

To address the limitations inherent in LSTM, Extended Long Short-Term Memory (xLSTM) proposes two significant alterations to the foundational LSTM concept outlined in Equation~\eqref{eq:lstm_idea}. The enhancements---namely, exponential gating and innovative memory architectures---broaden the LSTM framework by adding two novel variants: (i) the sLSTM (refer to Section~\ref{sec:sLSTM}), which employs a scalar memory, scalar update mechanism, and incorporates memory mixing; and (ii) the mLSTM (see Section~\ref{sec:mLSTM}), characterized by a matrix-based memory and an update rule utilizing covariances (outer products), which permits complete parallelization. Both variants incorporate exponential gating to augment standard LSTM capabilities. To facilitate parallel computation, mLSTM eliminates memory mixing, specifically the recurrent hidden-to-hidden connections. The architecture of both mLSTM and sLSTM supports the inclusion of multiple memory cells; sLSTM, in particular, enables memory mixing among these cells. Additionally, sLSTM supports a multi-head configuration where memory mixing occurs among cells within a head, but not across different heads. This adaptation, alongside exponential gating, introduces a novel form of memory interaction. In the context of mLSTM, the distinction between multiple heads and multiple cells becomes negligible, as the two configurations are effectively equivalent.

The incorporation of these novel LSTM variants within residual block structures yields xLSTM blocks, as discussed in Section~\ref{sec:xLSTMblock}. By organizing these xLSTM blocks in a residual manner, one constructs xLSTM architectures (refer to Section~\ref{sec:xLSTMarchitecure}). Figure~\ref{fig:xlstm_sketch} illustrates the overall xLSTM architecture and its constituent elements.

\subsection{Review of the Long Short-Term Memory}
\label{sec:orgLSTM}

The foundational concept behind LSTM~\citep{Hochreiter:91,Hochreiter:97nips,Hochreiter:97} centered on the development of a scalar memory cell, which serves both as the main site for computation and as a storage medium. This design effectively mitigates the vanishing gradient problem~\citep{Hochreiter:91,Hochreiter:00book} by utilizing the constant error carousel, achieved via cell state updates. Incorporated within the memory cell are three distinct gates: input, output, and forget gates, with the latter being introduced by \citet{Gers:00}. The set of equations governing the LSTM memory cell’s dynamics at time $t$ are:

\begin{align}
c_t \ &= \  f_t \ c_{t-1} \ + \ i_t \ z_t &  & &\text{cell state} \\
h_t  \ &= \ o_t \ \tilde{h}_t \ , & \tilde{h}_t \ &= \ \psi \left( c_t\right)
  &\text{hidden state} \\
z_t \ &= \ \varphi \left( \tilde{z}_t \right) \ , 
  &\tilde{z}_t \ &=  \ \mathbf{w}^\top_{z} \ \mathbf{x}_t \ + \
  r_{z}  h_{t-1} \ + \  b_{z} \ \
  &\text{cell input} \\
i_t \ &= \ \sigma \left( \tilde{i}_t  \right) \ , 
  &\tilde{i}_t \ &= \ \mathbf{w}^\top_{i} \ \mathbf{x}_t \ + \
  r_{i}  \ h_{t-1} \ + \  b_{i} \ \
  &\text{input gate} \\
f_t \ &= \ \sigma \left(  \tilde{f}_t \right) \ , 
  &\tilde{f}_t \ &= \ \mathbf{w}^\top_{f} \ \mathbf{x}_t  \ + \
  r_{f}  \ h_{t-1} \ + \  b_{f} \ \
  &\text{forget gate} \\
o_t \ &= \ \sigma \left( \tilde{o}_t \right) \ , 
  &\tilde{o}_t  \ &= \ \mathbf{w}^\top_{o} \ \mathbf{x}_t \ + \
  r_{o}  \ h_{t-1} \ + \  b_{o} \ \
  &\text{output gate} 
\end{align}

The vectors $\mathbf{w}_{z}$, $\mathbf{w}_{i}$, $\mathbf{w}_{f}$, and $\mathbf{w}_{o}$ serve as the input weight vectors linking the inputs $\mathbf{x}_t$ to the cell input, input gate, forget gate, and output gate, respectively. Analogously, $r_{z}$, $r_{i}$, $r_{f}$, and $r_{o}$ denote the recurrent weights that connect the previous hidden state $h_{t-1}$ with the cell input, input gate, forget gate, and output gate, in that order. The biases $b_{z}$, $b_{i}$, $b_{f}$, and $b_{o}$ are the associated offset parameters. The cell input and hidden state utilize the activation functions $\varphi$ and $\psi$ (commonly instantiated as $\tanh$). The function $\psi$ serves to constrain or squash the potentially unbounded values of the cell state. Every gate employs the sigmoid activation function, i.e., $\sigma \left( x \right)= 1/(1+\exp(-x))$. In subsequent developments, several scalar memory cells $\mathbf{c}_t \in \mathbb{R}^{d}$ were grouped into a vector, enabling the use of recurrent weight matrices $\mathbf{R} \in \mathbb{R}^{d \times d}$ for mixing outputs across memory cells~\citep{Greff:15}; see Appendix~\ref{sec:apporgLSTM} for a comprehensive discussion. Ablation experiments have demonstrated that each part of the memory cell plays an essential role~\citep{Greff:15}.

\subsection{sLSTM}
\label{sec:sLSTM}

To enable LSTMs to modify their storage choices dynamically, we incorporate exponential gates (represented in red), alongside mechanisms for normalization and stabilization. Specifically, both the input and forget gates may utilize exponential activation functions. For normalization purposes, we propose a normalizer state, which accumulates the sum of the input gate multiplied by the sequence of future forget gates.

The scalar sLSTM forward pass is given by:

\begin{align}
\label{eq:slstmforward}
c_t \ &= \  f_t \ c_{t-1} \ + \ i_t \ z_t & & 
 &\text{cell state} \\
n_t \ &= \  f_t \ n_{t-1} \ + \ i_t  & & 
 &\text{normalizer state} \\
h_t \ &= \ o_t \ \tilde{h}_t \ ,
  &\tilde{h}_t \ &= \ c_t / n_t
&\text{hidden state} \\
z_t \ &= \ \varphi \left( \tilde{z}_t \right) \ , 
  &\tilde{z}_t \ &=  \ \mathbf{w}^\top_{z} \ \mathbf{x}_t \ + \
  r_{z}  h_{t-1} \ + \  b_{z}
  &\text{cell input} \\
i_t \ &= \ \!\exp \left( \tilde{i}_t  \right) \ , 
  &\tilde{i}_t \ &= \ \mathbf{w}^\top_{i} \ \mathbf{x}_t \ + \
  r_{i}  \ h_{t-1} \ + \  b_{i}
  &\text{input gate} \\
f_t \ &= \ \sigma \left(  \tilde{f}_t \right) \ \text{OR} \
\!\exp \left(  \tilde{f}_t \right) \ ,
  &\tilde{f}_t \ &= \ \mathbf{w}^\top_{f} \ \mathbf{x}_t  \ + \
  r_{f}  \ h_{t-1} \ + \  b_{f}
  &\text{forget gate} \\
o_t \ &= \ \sigma \left( \tilde{o}_t \right) \ , 
  &\tilde{o}_t  \ &= \ \mathbf{w}^\top_{o} \ \mathbf{x}_t \ + \
  r_{o}  \ h_{t-1} \ + \  b_{o}
  &\text{output gate} 
\end{align}

The original LSTM gating mechanisms—specifically, those relying on the input and/or hidden state along with a bias term—are incorporated into the novel architectures. Since exponential activation functions have the potential to produce excessively large outputs that may result in overflow issues, we employ an auxiliary state \mbox{$m_t$~\citep{Milakov:18arxiv}} to enhance the stability of the gates:

\begin{align}
\label{eq:slstmstabil}
m_t \ &= \ \max \left( \log ( f_t ) + m_{t-1} , \log ( i_t ) \right) &\text{stabilizer state} \\
i'_t \ &= \ \!\exp \left( \log \left ( i_t \right) - m_t \right) \ = \ \!\exp \left( \tilde{i}_t - m_t \right) \ \, 
  &\text{stabil. input gate} \\
f'_t \ &= \ \!\exp \left( \log \left( f_t \right) + m_{t-1} - m_t \right) \ \, 
  &\text{stabil. forget gate}
\end{align}

As demonstrated in Appendix~\ref{sec:appsLSTM}, substituting $f_t$ with $f'_t$ and $i_t$ with $i'_t$ during the forward pass leaves both the overall network output and the gradients of the loss function with respect to the model parameters unaffected.

\paragraph{New Memory Mixing.} The sLSTM architecture, similar to the standard LSTM, is capable of accommodating several memory cells (refer to Appendix~\ref{sec:appsLSTM}). The presence of multiple memory cells permits mixing of memory through recurrent connections $\mathbf{R}_{\mathbf{z}}$, $\mathbf{R}_{\mathbf{i}}$, $\mathbf{R}_{\mathbf{f}}$, and $\mathbf{R}_{\mathbf{o}}$, which extend from the hidden state $\mathbf{h}$ to the memory cell input $\mathbf{z}$ as well as to the gates $\mathbf{i}$, $\mathbf{f}$, and $\mathbf{o}$. What differentiates the memory mixing in this approach is the utilization of exponential gating. In the new sLSTM, memory mixing can occur among memory cells within the same head, but not between different heads, as the model can employ multiple heads. The combination of head-based partitioning in sLSTM and exponential gating mechanisms offers a novel strategy for mixing memory.

\subsection{mLSTM}
\label{sec:mLSTM}

In order to boost the storage capabilities of LSTMs, we generalize the conventional scalar memory cell $c \in \mathbb{R}$ to a higher-dimensional matrix form, $\mathbf{C} \in \mathbb{R}^{d \times d}$. Consequently, the retrieval operation utilizes matrix multiplication. Specifically, at each time step $t$, the model encodes and stores a pair of $d$-dimensional vectors, namely the key $\mathbf{k}_t \in \mathbb{R}^d$ and the value $\mathbf{v}_t \in \mathbb{R}^d$—terms borrowed from the Transformer framework. Subsequently, at a later time $t + \tau$, the stored value $\mathbf{v}_t$ is intended to be accessed using a query vector $\mathbf{q}_{t+\tau} \in \mathbb{R}^d$. This configuration aligns with the framework of Bidirectional Associative Memories (BAMs) as established in \citep{Kohonen:72,Anderson:72,Nakano:72,Anderson:77}.

The process for incorporating a new key-value pair into memory is governed by the covariance update rule~\citep{Sejnowski:77,Dayan:91}, given by

\begin{align}
\mathbf{C}_t \ &= \ \mathbf{C}_{t-1} \ + \ \mathbf{v}_{t} \ \mathbf{k}_t^\top \ .
\end{align}

We presuppose the application of layer normalization prior to projecting inputs into keys and values, resulting in these representations possessing a mean of zero. As demonstrated by~\citet{Dayan:91}, the covariance update rule achieves optimality in terms of maximizing the separability among retrieved binary vectors, corresponding to the highest attainable signal-to-noise ratio. If one restricts retrieval operations to pairwise interactions and accepts the quadratic computational expense characteristic of attention mechanisms~\citep{Krotov:16,Krotov:17,Ramsauer:21}, an even greater degree of separability can be realized. Furthermore, the covariance update procedure is mathematically equivalent to that of Fast Weight Programmers~\citep{Schmidhuber:92ncfastweights,Schlag:21}, which have subsequently incorporated a fixed decay factor applied to $\mathbf{C}_{t-1}$ and a fixed learning rate applied to 
$\mathbf{v}_{t} \mathbf{k}_t ^\top$~\citep{Ba:16}.
Accordingly, we embed the covariance update scheme within the LSTM architecture, mapping the forget gate to the decay rate, the input gate to the learning rate, and employing the output gate as a modulator for the output vector.

Within this matrix memory framework, the normalizer state is constructed as a weighted aggregation of key vectors, with each vector scaled by the product of its corresponding input gate and the sequence of subsequent forget gates. This normalizer state effectively tracks the cumulative influence of the gating mechanisms. Notably, since the dot product of the query vector with the normalizer state may approach zero, we take the absolute value of this product and enforce a minimum threshold—commonly set at 1.0—following established practice~\citep{Sun:23arxiv}. The forward computation in mLSTM proceeds as follows:

\begin{align}
\label{eq:mlstm_recurrent_begin}
\mathbf{C}_t \ &= \  f_t \ \mathbf{C}_{t-1} \ + \ 
  i_t \ \mathbf{v}_t \ \mathbf{k}_t^\top & &  &\text{cell state} \\
\mathbf{n}_t \ &= \  f_t \ \mathbf{n}_{t-1} \ + \ 
  i_t \ \mathbf{k}_t & &  &\text{normalizer state} \\
\mathbf{h}_t  \ &= \ \mathbf{o}_t \ \odot \ \tilde{\mathbf{h}}_t
\ , \qquad \qquad 
\tilde{\mathbf{h}}_t \ = \   \mathbf{C}_t \mathbf{q}_t \ / \ 
  \max \left\{ |\mathbf{n}_t^\top \mathbf{q}_t|, 1 \right\} 
   &&&\text{hidden state} \\
\mathbf{q}_t \ &= \ \mathbf{W}_q \ \mathbf{x}_t \ + \ \mathbf{b}_q  & &  &\text{query input} \\
\mathbf{k}_t \ &= \ \frac{1}{\sqrt{d}} \mathbf{W}_k \ \mathbf{x}_t \ + \ \mathbf{b}_k  & &  &\text{key input} \\
\mathbf{v}_t \ &= \ \mathbf{W}_v \ \mathbf{x}_t \ + \ \mathbf{b}_v  & &  &\text{value input} \\
i_t \ &= \ \!\exp \!  \! \left( \tilde{i}_t  \right) \ , 
 \qquad \qquad \ \ \ \ \,
  \tilde{i}_t \ = \ \mathbf{w}^\top_{i} \ \mathbf{x}_t \ + \  b_{i} \ \ \ 
  &&&\text{input gate} \\
f_t \ &= \ \sigma \!  \left(  \tilde{f}_t \right) \ \text{OR} \ \!\exp \!  \! \left(  \tilde{f}_t \right) , \ \ \: \!
\tilde{f}_t \ = \ \mathbf{w}^\top_{f} \ \mathbf{x}_t  \ + \
  b_{f}
  &&& \text{forget gate} \\
\label{eq:mlstm_recurrent_end}
\mathbf{o}_t \ &= \ \sigma \left( \tilde{\mathbf{o}}_t \right) \ , \qquad \qquad \quad \ \ \ \,
  \tilde{\mathbf{o}}_t  \ = \ \mathbf{W}_{\mathbf{o}} \ \mathbf{x}_t \ + \
  \mathbf{b}_{\mathbf{o}}
  &&&\text{output gate} 
\end{align}

mLSTM, similar to the conventional LSTM, is capable of supporting several memory cells. In the context of mLSTM, having multiple heads is functionally the same as employing multiple cells, given the absence of memory mixing between them. To ensure stability for the exponential gating mechanisms in mLSTM, we adopt the same stabilization strategies used for sLSTM (refer to Equation~\ref{eq:slstmstabil}). Moreover, due to the lack of memory mixing in mLSTM, the associated recurrence relations can be reformulated to support parallel computation. Additional information is provided in Appendix~\ref{sec:appmLSTM}.

\subsection{xLSTM Architecture}
\label{sec:xLSTMarch}

\paragraph{xLSTM Blocks.}
\label{sec:xLSTMblock}
The purpose of an xLSTM block is to encode the sequence history through non-linear transformations in a high-dimensional feature space, facilitating the discrimination between distinct contexts or previous states. Effective separation of such historical information is essential for accurate next-element prediction within a sequence, for example, predicting subsequent tokens. This design principle draws upon Cover’s Theorem~\citep{Cover:65}, which posits that patterns which are not linearly separable in their original feature space are more likely to become linearly separable after being non-linearly projected into a higher-dimensional space. We examine two types of residual block designs: (i) A residual block incorporating a post up-projection step (analogous to the structure in Transformers), where the non-linear aggregation of prior states occurs within the initial dimensional space, followed by a linear up-projection into a space of higher dimensionality, the application of a non-linear activation, and a linear mapping that returns to the original space. This configuration is illustrated in the left side of Figure~\ref{fig:Backbones}
and depicted in the third column of Figure~\ref{fig:xlstm_sketch}. A comprehensive diagram of this version is provided in Figure~\ref{fig:appsLSTM_detailed} in the appendix. (ii) A residual block that performs pre up-projection (similar to architectures in State Space Models), where the mapping to a high-dimensional space is applied at the outset,

The high-dimensional space enables a nonlinear aggregation of historical information, which is subsequently mapped back to the original space via a linear transformation. In the case where an xLSTM block incorporates an sLSTM, we apply the post up-projection block. Conversely, for an xLSTM block integrating an mLSTM, the pre up-projection block is used because the memory capacity is increased in the expanded high-dimensional space. For further illustration, see the left side of Figure~\ref{fig:Backbones}, the third column in Figure~\ref{fig:xlstm_sketch}, and also Figure~\ref{fig:appsLSTM_detailed} in the appendix.

\paragraph{xLSTM Architecture. }
\label{sec:xLSTMarchitecure}
The xLSTM model is formed by sequentially layering individual blocks in a residual fashion~\citep{Srivastava:15,He:16}. Our design follows the prevalent approach in modern Large Language Models, utilizing pre-LayerNorm~\citep{Ba:16b} residual frameworks. For a visual summary, refer to the rightmost column of Figure~\ref{fig:xlstm_sketch}.

\subsection{Memory and Speed Considerations}

In contrast to Transformers, xLSTM networks feature linear computational requirements and maintain constant memory complexity relative to the length of the input sequence. Due to their compressive memory structure, xLSTM architectures are particularly advantageous for use in industrial contexts and for deployment on edge devices.

mLSTM's memory structure operates without the need for additional parameters; however, it incurs a high computational cost due to both the $d \times d$ memory matrix and the associated $d \times d$ update operations. Our approach balances the memory storage capacity with the complexity of computation. Despite the increased computational demand, these operations are well-suited to parallel execution on GPUs, resulting in only a minimal impact on the actual runtime.

Although mLSTM, like FlashAttention~\citep{Dao:22,Dao:23} or GLA~\citep{Yang:23arxiv}, supports parallel computation, sLSTM cannot be parallelized owing to the hidden-to-hidden memory mixing connections. Nonetheless, we introduced an efficient CUDA-based solution incorporating GPU memory optimizations down to the register level, resulting in performance that is usually within a factor of two of mLSTM’s speed.

\section{Related Work}
\label{sec:related}

\paragraph{Linear Attention.} A variety of approaches have been proposed to address the Transformer’s quadratic computational cost with respect to context length, aiming to achieve linear attention mechanisms. The Synthesizer constructs attention weights in a synthetic manner, eliminating explicit token-to-token dependencies~\citep{Tay:20arxiv}. Linformer achieves self-attention through a low-rank factorization of the attention matrix, allowing for a linear approximation~\citep{Wang:20arxiv}. The Linear Transformer applies a linearization to the attention computation itself~\citep{Katharopoulos:20}. The Performer utilizes a positive orthogonal random features strategy to provide a linear approximation of the softmax function in attention~\citep{Choromanski:21}. In contrast, Structured Global Convolution (SGConv)~\citep{Li:22} and the Hyena Hierarchy~\citep{Poli:23} substitute the attention module with highly efficient long-range convolutional operations.

\paragraph{State Space Models.} In recent developments, State Space Models (SSMs) have garnered considerable attention due to their linear computational complexity relative to sequence length and their competitive results against Transformers. Early notable work includes the Structured State Space sequence model (S4)~\citep{Gu:21}, which was soon followed by models such as the Diagonal State Space (DSS) model~\citep{Gupta:22}, Gated State Space (GSS) models~\citep{Mehta:22}, S5 model~\citep{Smith:22}, Bidirectional Gated SSM (BiGS)~\citep{Wang:22}, H3 model~\citep{Fu:23}, as well as the Mamba architecture~\citep{Gu:24arxiv}.

\paragraph{Recurrent Neural Networks.} Recurrent Neural Networks (RNNs) have emerged as candidates to supplant Transformers and attention mechanisms, largely owing to their linear scaling with sequence length. Notable advances include RNNs equipped with Deep Linear Recurrent Units (LRUs), which have demonstrated encouraging performance on language modeling benchmarks~\citep{Orvieto:23, De:24arxiv}. Similarly, the Hierarchically Gated Linear RNN (HGRN)~\citep{Qin:23} and its successor HGRN2~\citep{Qin:24arxiv} have yielded favorable outcomes. Among RNN-based large language models, the RWKV framework~\citep{Peng:23arxivshort,Peng:24arxivshort} stands out for achieving results that are comparable with those of Transformer-based architectures.

\paragraph{Gating.}
Gating represents a fundamental concept introduced by LSTM, and has subsequently been re-envisioned and integrated into a broad range of modern techniques. Various models have incorporated gating mechanisms, including HGRN~\citep{Qin:23}, HGRN2~\citep{Qin:24arxiv}, Gated Linear Attention (GLA)~\citep{Yang:23arxiv}, Gated State Space (GSS) models~\citep{Mehta:22}, Bidirectional Gated SSM (BiGS)~\citep{Wang:22}, Moving Average Equipped Gated Attention (MEGA)~\citep{Ma:22}, RWKV~\citep{Peng:23arxivshort}, and Mamba~\citep{Gu:24arxiv}.

\paragraph{Covariance Update Rule.} In order to improve storage capacity, we integrated a matrix memory governed by a covariance update rule into the mLSTM cell. Comparable strategies employing this kind of update mechanism can be found in Fast Weight Programmers \citep{Schmidhuber:92ncfastweights,Schlag:21}, RWKV-5 and RWKV-6~\citep{Peng:24arxivshort}, Retention~\citep{Sun:23arxiv}, Linear Transformer~\citep{Katharopoulos:20}, and HGRN2~\citep{Qin:24arxiv}.

\paragraph{Most Related.} The models most akin to xLSTM, from a conceptual standpoint, are Retention~\citep{Sun:23arxiv}, RWKV~\citep{Peng:23arxivshort,Peng:24arxivshort}, and HGRN2~\citep{Qin:24arxiv}. These architectures incorporate mechanisms such as matrix memory and/or gating. Nonetheless, unlike the new sLSTM, these methodologies do not support memory mixing. The capability for memory mixing is crucial for addressing state tracking challenges, making LSTMs inherently more expressive compared to State Space Models (SSMs) and Transformers~\citep{Merrill:24,Deletang:23}. Effective state tracking is essential for tasks like program evaluation and maintaining entity information throughout extended narratives.

\paragraph{Residually Stacking Architectures.} In alignment with the design paradigm of nearly all modern large-scale deep neural networks, xLSTM architectures are realized by stacking component modules with residual connections~\citep{Srivastava:15,He:16}.

This architectural approach was pivotal for the advancement of very deep convolutional networks~\citep{He:16} and the Transformer architecture~\citep{Vaswani:17}. Transformers, in particular, form the backbone of Large Language Models (LLMs) such as GPT-3~\citep{Brown:20short}, ChatGPT~\citep{Schulman:22short}, GPT-4~\citep{Achiam:23gpt4short}, Megatron-LM~\citep{Shoeybi:19}, Gopher~\citep{Rae:21short}, ERNIE 3.0 Titan~\citep{Wang:21arxivshort}, GLaM~\citep{Du:21glamshort}, Chinese M6~\citep{Lin:21}, multilingual AlexaTM 20B~\citep{Soltan:22}, OPT~\citep{Zhang:22opt}, Chinchilla~\citep{Hoffmann:22short}, BLOOM~\citep{Scao:22arxivshort}, GLM-130B~\citep{Zeng:22arxivshort}, LaMDA~\citep{Thoppilan:22short}, PaLM~\citep{Chowdhery:22short}, Llama~\citep{Touvron:23arxiv}, and Gemini~\citep{GeminiTeam:23arxiv,Reid:24arxiv}.

\section{Experiments}
\label{sec:Exp}

In this section, we present an empirical evaluation of xLSTM, benchmarking its performance against current approaches in the area of language modeling. Section~\ref{sec:ExpSyntethic} explores the particular strengths of xLSTM on a suite of synthetic tasks. Subsequently, in Section~\ref{sec:ExpComparison}, we report validation set perplexity for a range of contemporary language modeling architectures, all trained on a 15B token sample from SlimPajama~\citep{Soboleva:23}. We also carry out ablation analyses on xLSTM under these training conditions. Furthermore, we study how the various models scale, adopting the methodologies used by \citet{Kaplan:20} and \citet{Brown:20short}.

Section~\ref{sec:ExpLanguage} is dedicated to an extensive language modeling comparison. Here, both xLSTM and the leading architectures identified in Section~\ref{sec:ExpComparison} are trained with 300B tokens from SlimPajama~\citep{Soboleva:23}. Our evaluation proceeds in several steps: first, we examine the models’ abilities to generalize to extended sequence lengths; next, we evaluate each approach in terms of validation perplexity as well as downstream task effectiveness~\citep{Sutawika:24}; third, we measure performance across the 571 textual domains of the PALOMA language benchmark~\citep{Magnusson:23arxivshort}; and fourth, we revisit the scaling analysis for each method, now leveraging a dataset that is 20 times larger than in our earlier experiments.

\label{sec:xlstm_blocks_notation}
Throughout our experiments, we denote xLSTM[$a$:$b$] to represent a configuration where the proportion of mLSTM-based blocks to sLSTM-based blocks is $a/b$. As an illustration, xLSTM[7:1] corresponds to a composition of eight blocks, of which seven are mLSTM-based and one is sLSTM-based. For a standard total of 48 blocks, this configuration would comprise 42 mLSTM-based and 6 sLSTM-based blocks. Additionally, in all experimental setups, mLSTM blocks utilize pre up-projection, while sLSTM blocks employ post up-projection.

\subsection{Synthetic Tasks and Long Range Arena}
\label{sec:ExpSyntethic}
We begin by evaluating the capability of xLSTM’s exponential gating combined with memory mixing on formal language benchmarks~\citep{Deletang:23}. Subsequently, we investigate how xLSTM’s matrix-based memory performs in the Multi-Query Associative Recall task~\citep{Arora:23arxiv}. Lastly, we examine xLSTM’s ability to handle extended sequence lengths using the Long Range Arena suite~\citep{Tay:21}.

\paragraph{Evaluation of xLSTM’s Exponential Gating and Memory Mixing Capabilities.} We assess the newly introduced exponential gating mechanism with memory mixing in xLSTM, designed to tackle state tracking challenges~\citep{Merrill:24, Merrill:23}. Our experiments build upon and extend the formal language benchmarks from \citet{Deletang:23} to allow for training over sequences of multiple lengths, facilitating length extrapolation. Additional information on the experimental setup and expanded results are provided in Appendix~\ref{sec:appExpSynthetic-formalLang}. 

xLSTM’s performance is benchmarked against a range of architectures including Transformers, State Space Models, and traditional Recurrent Neural Networks. We report accuracy specifically over task-relevant tokens, normalizing performance between 0 (chance) and 1 (optimal). Comparisons are performed on 2-block variants of the following methods: xLSTM[0:1] (equivalent to sLSTM), xLSTM[1:0] (equivalent to mLSTM), xLSTM[1:1], Llama, Mamba, RWKV, Retention, Hyena, LSTM, and LSTM implemented within Transformer blocks (LSTM (Block)). Outcomes from these experiments are presented in Figure~\ref{fig:formal-main}.

Notably, models such as Transformers and State Space Models that lack memory mixing—and thus effective state tracking—are unable to solve tasks involving regular grammars, such as the parity task. This observation aligns with prior findings indicating that Transformers and State Space Models possess fundamentally weaker computational power compared to RNNs~\citep{Merrill:24, Merrill:23, Deletang:23}.

\paragraph{Test of xLSTM's Memory Capacities on Associative Recall Tasks.} This study investigates the capacity of xLSTM's newly introduced matrix memory by employing the Multi-Query Associative Recall task~\citep{Arora:23arxiv}. In this task, sequences are constructed by randomly assigning key-value pairs drawn from an extensive vocabulary; the model must memorize these associations for subsequent recall. To increase the challenge relative to previous versions, we scale the number of key-value pairs up to 256 and permit context lengths as long as 2048, enabling a comprehensive assessment of model memory capabilities. Our comparison spans 2-block versions of Llama, Mamba, RWKV-5, RWKV-6, xLSTM[1:1], and xLSTM[1:0], with performance measured by recall accuracy. Due to their memory scaling exponentially with the coding dimension~\citep{Ramsauer:21}, Transformers such as Llama represent the reference point for this benchmark. Figure~\ref{fig:mqar-main} presents the outcomes. Among non-Transformer approaches, xLSTM[1:1] achieves superior results, maintaining this advantage even at smaller model scales. Notably, the inclusion of the sLSTM block does not reduce memory capacity; on the contrary, it enhances performance, particularly in the hardest scenario involving 256 pairs. Further analysis, including extrapolation to longer contexts than those used in training, is detailed in Appendix~\ref{sec:appExpSynthetic-mqar}, highlighting xLSTM's expanded memory proficiency.

\paragraph{Test of xLSTM's Long Context Capabilities on Long Range Arena.} To evaluate how well xLSTM handles lengthy sequences and extended contexts, we benchmark various approaches using the Long Range Arena~\citep{Tay:21}. The xLSTM consistently achieves high performance across all tasks, indicating its architectural effectiveness and robustness for a range of long context challenges.

Additional information is provided in Appendix~\ref{sec:appExpSynthetic-lra}.

\subsection{Method Comparison and Ablation Study}
\label{sec:ExpComparison}

In order to investigate the central question posed by this study—namely, the capabilities of our proposed LSTM variants when scaled for language modeling—we conduct experiments using xLSTMs, Transformers, State Space Models, and additional approaches. Each model is trained in an auto-regressive fashion on 15B tokens from SlimPajama. Subsequently, we evaluate model performance on the validation dataset and carry out ablation analyses specifically focused on the xLSTM architectures.

\paragraph{Comparing xLSTM to Other Methods.} For the purposes of benchmarking, all models are trained on a dataset comprising 15B tokens sourced from SlimPajama~\citep{Soboleva:23}. Model evaluation is conducted using validation set perplexity. The following models are included in our comparison: xLSTM (introduced in this work), GPT-3 (Transformer)~\citep{Brown:20short}, Llama (Transformer)~\citep{Touvron:23arxiv}, H3 (SSM)~\citep{Fu:23}, Mamba (SSM)~\citep{Gu:24arxiv}, RWKV-4 (RNN)~\citep{Peng:23arxivshort}, RWKV-5 (RNN)~\citep{Peng:24arxivshort}, RWKV-6 (RNN)~\citep{Peng:24arxivshort}, GLA (linear Transformer)~\citep{Yang:23arxiv}, HGRN (RNN)~\citep{Qin:23}, HGRN2 (RNN)~\citep{Qin:24arxiv}, RetNet (linear Transformer)~\citep{Sun:23arxiv}, Hyena (linear Transformer)~\citep{Poli:23}, xLSTM[1:0], and xLSTM[7:1] (see Section~\ref{sec:xlstm_blocks_notation}). 
Models were trained under mixed-precision regimes; for RWKV-5, RWKV-6, GLA, and HGRN2, the mixed-precision implementation did not rely on PyTorch’s automatic mixed precision framework (refer also to Appendix Section~\ref{sec:appGenTrainProc}).
We organize the models into three primary groups: (a) Transformers, (b) State Space Models (SSMs), and (c) Recurrent Neural Networks (RNNs) in conjunction with linear Transformers. Notably, linear Transformers implement a linear alternative to traditional attention. All models are designed to match the parameter count of a GPT-3 variant with 350M parameters, specifically with embedding dimension set to 1024 and comprising 24 residual layers. Of these, only GPT-3 employs tied weights for token and output embeddings, thereby resulting in a smaller parameter count. As displayed in Table~\ref{tab:spaj15b_model_results}, xLSTM achieves superior validation perplexity compared to all other evaluated models. Further experimental details are discussed in Appendix~\ref{sec:appExpComparison}. 
The scaling trends for this experimental setting are depicted in Figure~\ref{fig:temp_sclaw15B}, suggesting that xLSTM maintains its performance advantages as model size increases.

\paragraph{Ablation Studies.}
Table~\ref{tab:spaj15b_model_results} alongside Figure~\ref{fig:temp_sclaw15B} illustrate that xLSTM attains strong language modeling performance when trained with 15B tokens from the SlimPajama dataset. In order to isolate the contributions that differentiate xLSTM from a standard LSTM, we incrementally modify a basic LSTM architecture towards xLSTM. The process begins with the insertion of LSTM layers within pre-LayerNorm residual frameworks. Next, the design is extended to include a block featuring post up-projection. The final step introduces exponential gating mechanisms and a matrix memory component.

Table~\ref{tab:ablstudies} (top) presents the experimental findings.

Through ablation experiments, we determine that the significant performance gains arise from the combination of exponential gating and matrix memory. Given the established relevance of gating in both RNNs and State Space Models, we further examine the impact of various gating configurations. The results in Table~\ref{tab:ablstudies} (bottom) indicate that enabling each gate to be trainable and responsive to input leads to consistent improvements. Further analysis of specific backbone components is provided in Appendix~\ref{sec:appAblDetails}.

\subsection{xLSTM as Large Language Model}
\label{sec:ExpLanguage}

This work concludes with extensive language modeling experiments at scale, where we evaluate xLSTM in the context of large language models. To this end, we expand the training dataset, utilizing 300B tokens sourced from SlimPajama. This training size aligns with the token count employed by models such as Mamba~\citep{Gu:24arxiv} and Griffin~\citep{De:24arxiv}.

We evaluate xLSTM in comparison to RWKV-4, Llama, and Mamba, each representing a distinct model class as described in Section~\ref{sec:ExpComparison}. RWKV-4 is chosen to represent RNNs because appropriate training precision configurations for RWKV-5, RWKV-6, and HGRN2 were only established after commencing training on the 300B token experiments (refer to Appendix~\ref{sec:appGenTrainProc} for details). Our experiments include models of varying sizes (125M, 350M, 760M, 1.3B), and all variants are evaluated for their ability to extrapolate to longer sequences as well as their validation set performance. In addition, we examine performance on a range of downstream tasks, conduct language modeling assessments across 571 text domains from the PALOMA benchmark, and analyze the models’ scaling behavior.

\paragraph{Sequence Length Extrapolation.}
\label{sec:spaj300B_ctx_extrapolation}
We begin by evaluating the sequence length extrapolation capabilities of 1.3B-parameter models across xLSTM, RWKV-4, Llama, and Mamba architectures. Each model undergoes training with a fixed context length of 2048 tokens and is subsequently assessed on extended context lengths reaching up to 16384 tokens. The outcomes are presented in Figure~\ref{fig:spaj300B_seq_extrapolation}. Notably, xLSTM outperforms the other architectures by sustaining low perplexity as sequence lengths increase.

\paragraph{Validation Perplexity and Downstream Tasks.}
\label{sec:spaj_downstream_eval}
Next, across each model size, we assess xLSTM, RWKV-4, Llama, and Mamba models by measuring their next token prediction performance on the SlimPajama validation set. Additionally, we test these models on downstream benchmarks that evaluate their capabilities in common sense reasoning.

The third column in Table~\ref{tab:lmeval} presents the validation set perplexity values for the various approaches. Across every model size, xLSTM[1:0] and xLSTM[7:1] consistently achieve the lowest perplexity on the validation set. The remaining columns of Table~\ref{tab:lmeval} report results on downstream evaluations. In most tasks and for all model scales, xLSTM outperforms competing methods; an exception is the ARC benchmark, where Mamba occasionally obtains the highest score. Additional results can be found in Appendix~\ref{sec:appExpLanguage}.

\paragraph{Performance on PALOMA Language Tasks.} Third, we evaluate the next token prediction accuracy of the xLSTM, RWKV-4, Llama, and Mamba architectures across all model capacities using the PALOMA language tasks~\citep{Magnusson:23arxivshort}. We report results in terms of perplexity for next token prediction, spanning 571 distinct text domains—including sources such as nytimes.com and the r/depression subreddit. In Table~\ref{tab:ppleval}, we present the token prediction perplexity, categorizing results by language modeling tasks (first seven columns) and specific domain evaluations (last five columns). The xLSTM[1:0] variant consistently outperforms xLSTM[7:1] in these language modeling assessments.

xLSTM[1:0] achieves lower perplexity than Mamba in 568 out of 571 text domains (99.5\%), outperforms Llama in 486 of 571 cases (85.1\%), and surpasses RWKV-4 in 570 of 571 domains (99.8\%). Further information can be found in Appendix~\ref{sec:appExpLanguage}.

\paragraph{Scaling Laws.} Fourth, we investigate power-law scaling trends, which provide a basis for predicting model performance as size increases~\citep{Kaplan:20,Brown:20short}. The scaling characteristics are depicted in Figure~\ref{fig:temp_sclaw300B}. While all evaluated models exhibit comparable scaling trends, their starting points differ. RWKV-4 demonstrates the lowest performance, with Llama and Mamba ranking above it in succession. xLSTM outperforms Mamba by roughly the same margin that separates Mamba from Llama. The observed scaling suggests that as model size grows, xLSTM is likely to maintain a performance advantage over both Transformer-based and State-Space architectures.

\textbf{Generation Times and Maximal Throughput.} To assess performance, we report the generation latency in Figure~\ref{fig:inference_time_generation} and evaluate peak decoding throughput in Figure~\ref{fig:inference_time_decoding_throughput} (left) for xLSTM models with 1.3B parameters. Comparative analysis includes analogous Mamba, Llama, and RWKV models sourced from HuggingFace, with Llama augmented by a static key-value cache. During our investigations, compiling the Transformer Model with a full cache and leveraging \texttt{torch.compile} for Mamba were both unsupported. In text generation benchmarks, we uniformly adopt a batch size of 1 and set the pre-fill length to 16 for all models—a scenario that particularly benefits Transformer-based architectures. 

As depicted in Figure~\ref{fig:inference_time_generation}, the xLSTM, alongside recurrent approaches such as Mamba and RWKV-4, demonstrates a linear increase in generation time, whereas Llama’s computation grows quadratically with sequence length. For throughput evaluations, we vary batch size and prefill settings for Llama to test its limits. Figure~\ref{fig:inference_time_decoding_throughput} (right) further reveals that xLSTM is capable of handling substantially larger batches than Llama because of its fixed memory usage, resulting in superior throughput across tested configurations.

\section{Limitations}

(i) Unlike mLSTM, the memory mixing mechanism in sLSTM inherently restricts parallel operations, thereby preventing efficient parallel implementation. Despite this, we have engineered a high-speed CUDA kernel for sLSTM; this kernel currently operates at less than double the runtime of our parallelized mLSTM kernel. (ii) The existing CUDA kernels for mLSTM lack optimization, resulting in execution that is approximately four times slower than the performance of FlashAttention or the scan algorithm adopted in Mamba. The development of more efficient CUDA kernels, inspired by approaches like FlashAttention, has the potential to significantly accelerate mLSTM. (iii) The matrix memory structure in mLSTM introduces substantial computational demands, as it necessitates handling $d \times d$ matrices. However, since both the memory update and retrieval stages are parameter-free and readily lend themselves to parallelization through conventional matrix operations, the practical increase in wall clock time attributable to the matrix memory’s complexity remains minimal. (iv) It is crucial to carefully select the initialization strategy for the forget gates. (v) As the matrix memory operates independently of the sequence length, scaling up the sequence length can lead to memory saturation when dealing with extended context windows. Nevertheless, empirical observations suggest this issue does not emerge for context lengths up to 16k; refer to Section~\ref{sec:spaj300B_ctx_extrapolation} for further details. (vi) Given the intensive computational requirements associated with large-scale language modeling experiments, we have not undertaken full optimization of the architecture or hyperparameters, particularly regarding the larger xLSTM variants. We expect that substantial and thorough optimization will be necessary for xLSTM models to realize their full capabilities.

\section{Conclusion}
We have provided a partial response to our central question: What are the capabilities of language modeling when LSTM architectures are scaled to the range of billions of parameters? Based on our results, the answer is: ``At a minimum, they reach the levels achieved by prevailing methods such as Transformers and State Space Models.'' By introducing xLSTM—an advancement of LSTM featuring exponential gating with memory mixing and a revised memory architecture—we demonstrate that these models are competitive with, and often surpass, state-of-the-art techniques in language modeling. Our exploration of scaling laws further suggests that enlarging xLSTM models may position them as formidable alternatives to contemporary Large Language Models, which primarily utilize Transformer architectures. Furthermore, the innovations behind xLSTM could have significant implications for adjacent domains including Reinforcement Learning, Time Series Prediction, and the simulation of physical phenomena.

\end{document}